\section{Finite-order twisting preserves the same ordinary ideal}
\label{sec:hida-twisting-v45}

This chapter proves the literal $B.\mathrm{TwistOk}(J)$, in every
original positive rank.  It reuses HPT17's ordinary return theorem;
the rank-$2$--$4$ descent input in HPT18 is not used.  Fix all
original Hida, coefficient, level, type and Hecke data as in HPT,
with $V.S'=s.S$, $n\geq1$, $\ell$ odd and
$R_H.\mathrm{IsJ}(B,K,J)$.  The isomorphism $\iota$ remains fixed.
We retain the same original global Artin, local classification,
local correspondence and period/monodromy foundations as before.
Every occurrence of a local correspondence means the complete
original $\mathrm{IsRecT}$ quantifier; no new global descent
foundation is introduced.

The exact input consists of two original points $x,x'$, a weight
$h$, and a character $\beta:G_0\to\overline{\mathbb Q}_\ell^\times$
such that
\[
 \begin{gathered}
 J\subseteq\ker x,\quad x.\mathrm{Good}(h),\quad
 \ker\beta\text{ open},\quad \beta^{\ell^k}=1\text{ for some }k,\quad
 \beta\circ\mathrm{conj}=\beta^{-1},\\
 I_u\subseteq\ker\beta\quad(u\in s.S_t),\qquad
 I_0\subseteq\ker\beta,\qquad
 x.D.\mathrm{IsTwist}(\beta,x'.D),\\
 x'.\psi_{v,i}=x.\psi_{v,i}
                 \bigl(\beta\circ p|_{D_{s.A(v)}}\bigr)
                                        \quad\text{for all }v,i.
 \end{gathered}                                                  \tag{HTS.1}
\]
Here $p=V.\mathrm{pr}$, and all inertia groups are the original
frame groups.  Goodness includes arithmeticity and injectivity
of the ordered finite-part characters at each selected place.
The desired conclusion is $J\subseteq\ker x'$ on the same ring.

% BEGIN HTS1
\paragraph{HTS1. An odd-order conjugate-dual character supplies the required unitary character.}
\label{claim:HTS1}
Put $\mu=\beta\circ p$ and $m=\ell^k=2N+1$.  There is a
finite-order Hecke character $\nu$ with Galois character
$\delta=\mu^{N+1}$, and $\kappa=\nu^2$ has Galois character
$\mu$ with weight zero.  The character
$\chi=\nu\circ\det$ on the actual adelic unitary group has
open kernel and is trivial on rational and infinite elements.
At a split local factor its character is $\kappa_w\circ\det$;
at a nonsplit local factor it is $\nu_w\circ\det$.
\begin{proof}
An open kernel makes $\beta$ locally constant, hence continuous;
the same holds after composing with the continuous $p$.
Since $m$ is odd, $\delta^2=\mu^{m+1}=\mu$ and $\delta^m=1$.
It is conjugate-dual because $p$ carries the specified complex
conjugation to the original frame involution.

Choose the original global Artin map.  The continuous character
$\delta$ kills the closure of the commutator subgroup, so factors
through the abelianized Galois group.  Compose this factor with
Artin and $\iota$ to define $\nu$ on ideles.  Its image is finite,
so its kernel is open despite the absence of any continuity
assumption on $\iota$.  The Artin kernel contains rational
ideles, so $\nu$ is an actual Hecke character.  At weight zero
the defining $\mathrm{IsLAdicOf}$ equality is simply
\[
         \delta(g)=\iota^{-1}\nu(a)
                       \quad\text{if }\mathrm{Art}(a)=\bar g.
                                                               \tag{HTS.2}
\]
This proves the stated attachment, its order bound, and the
attachment of $\kappa$ to $\mu$.  Uniqueness follows by applying
(HTS.2) to every idele, since every element of the abelianized
Galois group has a lift to the Galois group.

The conjugate-dual identity gives $\nu(\bar a)=\nu(a)^{-1}$.
One may check this without presupposing functoriality of the
chosen Artin map: outside the union of the two finite conductor
sets, conjugating a Frobenius lift and its prime preserves its
residue power and gives the claimed inverse values.  Finite
characters of the idele class group correspond to continuous
finite characters of the Artin quotient.  Indeed their open
kernels contain rational ideles and the connected infinity
factor, hence the closure specifying the original Artin kernel;
the quotient map factors them.  Chebotarev equality at the above
cofinite Frobenius set then identifies these characters, proving
the identity on all ideles.

The determinant character $\chi$ has open kernel because the
determinant is continuous and $\nu$ has open kernel.  On a
rational matrix its determinant is a rational idele, so it is
trivial.  It is trivial on infinity: $L$ is totally complex,
and every infinity idele has an $m$-th root with finite part one;
applying the order-$m$ character gives value one.
At a split pair $w,\bar w$, the unitary equation says that the
conjugate of the determinant at $\bar w$ is the inverse of the
determinant at $w$.  The two local $\nu$ values are consequently
equal, and their product is $\nu_w^2=\kappa_w$.  At a nonsplit
place there is one coordinate, giving $\nu_w$ itself.
This proves the precise split/nonsplit local formula on the
actual adelic group, including the square in the split case.
\end{proof}
% END HTS1

% BEGIN HTS2
\paragraph{HTS2. Multiplication by the actual determinant character twists automorphic spaces.}
\label{claim:HTS2}
For an actual irreducible automorphic subspace $\sigma$, set
\[
          M_\chi f(g)=\chi(g)f(g),\qquad \sigma'=M_\chi\sigma.
                                                               \tag{HTS.3}
\]
Then $\sigma'$ is actual irreducible automorphic, with the same
infinity weights.  At every split place its actual components
are precisely the twists of those of $\sigma$ by
$\kappa_w\circ\det$.  Wherever $\mu$ is trivial on inertia,
unramifiedness is preserved; so is every fixed vector for a
subgroup of the local maximal compact, in particular $K_c(1)$.
\begin{proof}
$M_\chi$ is a linear isomorphism with inverse $M_{\chi^{-1}}$,
and direct substitution gives
\[
          R_U(g)M_\chi f=\chi(g)M_\chi R_U(g)f.                 \tag{HTS.4}
\]
Continuity and left rational invariance follow from HTS1.
Intersect an open stabilizer of $f$ with $\ker\chi$ to get
an open stabilizer of $M_\chi f$.  The infinity orbit span
is the image of the original finite-dimensional span, so it
is finite-dimensional.  Thus the original automorphic-form
conditions all hold.  Equation (HTS.4) preserves invariant
subspaces in both directions, proving irreducibility and
nonvanishing.  Triviality of $\chi$ at infinity preserves
every infinity intertwiner and hence its weight.

Compose any local-component embedding with $M_\chi$.
Equation (HTS.4) and the split formula in HTS1 give exactly
the twisted local representation.  Conversely use the inverse
map and inverse twist.  This proves the statement for every
actual component, not just a chosen family.  A finite-order
Galois character is trivial on inertia exactly when its
weight-zero Hecke character is trivial on local units, by
local Artin surjectivity.  If $\mu$ is trivial on inertia,
so is its power $\delta$; thus both $\kappa_w$ and $\nu_w$
are unramified.  The determinant of a maximal-compact matrix
is a unit, so the local character of HTS1 is one on the
compact subgroup in question.  The same nonzero vector
remains fixed under the twisted action, at split and nonsplit
places alike.
\end{proof}
% END HTS2

% BEGIN HTS3
\paragraph{HTS3. The two exact guards for local twisting.}
\label{claim:HTS3}
Two local situations provide precisely the original
$\mathrm{Good}(\tau)$ and $\mathrm{Gd}(\mathrm{rc}(\mathrm{algNorm}\tau))$
guards: first, a prescribed zero-monodromy inertia type given by
restrictions of global characters; second, an actual twist
$\tau\otimes(\kappa_w\circ\det)$ which is spherical.
These statements hold for every permitted correspondence and
do not assert $\mathrm{PlainTor}$ for a nongeneric component.
\begin{proof}
For the first situation, write the parameter using the original
Langlands multisegment law.  The type hypothesis gives zero
monodromy and the inertia polynomial of restrictions of global
characters.  For a Weil element $s$ and an inertia element $t$,
their commutator belongs to inertia and each such global character
is one on it.  Hence its total characteristic polynomial is
$(X-1)^n$.  The polynomial of any irreducible Weil constituent
of any supercuspidal support factor divides this polynomial.
Its inertia action is semisimple, so the commutator acts as the
identity.  Thus inertia commutes with the entire irreducible
Weil image and is scalar by Schur's lemma.  An eigenline for
Frobenius is then stable under inertia and Frobenius, which
generate the Weil group.  Irreducibility forces dimension one.
The rank equality of the correspondence makes every cuspidal
support factor rank one.  This is enough for the original
$\mathrm{Good}$ disjunction, including repeated or linked
characters.  The support product law makes the polynomial on
every commutator of two Weil elements equal to $(X-1)^n$;
the type already gives zero monodromy.  These are exactly
the two fields of $\mathrm{Gd}$.

For the second situation, GUA12 identifies the spherical
component's support as unramified characters and its parameter
as a sum of rank-one unramified characters with zero monodromy.
Untwisting its support by the single character $\kappa_w$
again gives rank-one support, hence $\mathrm{Good}(\tau)$.
The twist law of $\mathrm{IsRecT}$ applies because this guard
holds.  It identifies the untwisted parameter with the same
sum tensored by a rank-one Weil character.  Such a character
is one on every commutator and does not change whether the
monodromy is zero.  The untwisted parameter is therefore
$\mathrm{Gd}$ at every correspondence.  Neither argument
uses genericity or a stronger all-local compatibility theorem.
\end{proof}
% END HTS3

% BEGIN HTS4
\paragraph{HTS4. Finite-character twisting respects the full away local comparison.}
\label{claim:HTS4}
Suppose $w\nmid\ell$, $R$ is continuous, $\mu$ is the finite
character of HTS1, $\tau$ satisfies the original $\mathrm{Good}$
guard, and $\tau'\simeq\tau\otimes(\kappa_w\circ\det)$.
For every original $\mathrm{IsRecT}$,
\[
 \mathrm{IsWDOf}(\iota,R,\mathrm{rc}(\mathrm{algNorm}\tau))
 \ \Longrightarrow\
 \mathrm{IsWDOf}(\iota,R\mu,\mathrm{rc}(\mathrm{algNorm}\tau')).
                                                               \tag{HTS.5}
\]
Here $(R\mu)(g)=\mu(g)R(g)$ on the same vector space.
\begin{proof}
Choose a finite coefficient field containing the representation's
compact image and the finite image of $\mu$, as in ADK1.
Intersect the open subgroup used for the original away local
monodromy witness with $\ker\mu$.  On it $R\mu=R$, so the
same nilpotent monodromy operator works.  On the Weil group
the resulting raw action is the old action multiplied by
$\iota\mu(s)$.  This follows immediately from the defining
formula $R(s)\exp(-t_\ell(s)N)$.  A scalar factor commutes with
the action and its semisimplification correction; hence the
Frobenius-semisimplified parameter is the original one tensored
by this rank-one character.  The same coefficient lattice,
after the harmless finite enlargement, supplies an actual
original $\mathrm{IsWDOf}$ witness.

Local-global Artin compatibility and (HTS.2) identify
$\iota\mu(s)=\kappa_w(a(s))$.  OFL1 identifies this $a$ with
the map in the chosen correspondence.  Its twist law applies
under $\mathrm{Good}(\tau)$; the common algebraic normalizing
norm twist commutes with the character twist.  It gives exactly
the parameter on the right of (HTS.5).  The choices were made
for an arbitrary correspondence and prolongation, so all those
quantifiers remain intact.
\end{proof}
% END HTS4

% BEGIN HTS5
\paragraph{HTS5. Good finite parts give a full principal series and its exact twist.}
\label{claim:HTS5}
Suppose $\tau$ is a smooth irreducible subquotient of $\mathrm{PS}(\omega)$,
$\mathrm{OrdChar}(\omega_i,h_i,\psi_i)$ holds, and the finite
parts $\phi_i$ of these characters are pairwise distinct.
Then $\tau\simeq\mathrm{PS}(\omega)$.  If $\tau'$ is its twist
by $\kappa_w\circ\det$, then
\[
 \tau'\simeq\mathrm{PS}(\omega'),\qquad
 \omega'_i=\omega_i\kappa_w,\qquad
 \mathrm{OrdChar}(\omega'_i,h_i,\psi_i\mu|_{D_w}).               \tag{HTS.6}
\]
The latter characters are locally algebraic with the same weights.
\begin{proof}
By HPT3 the restrictions of $\omega_i$ to units are inverses,
through $\iota$, of the finite parts on Artin images of inertia.
Surjectivity of local Artin on units implies that the restrictions
are pairwise distinct.  No two $\omega_i$ therefore differ by
any unramified norm twist, including twists $-1,0,1$.
The original totally disjoint product theorem and induction
in stages make $\mathrm{PS}(\omega)$ irreducible.  Its nonzero
irreducible subquotient $\tau$ is isomorphic to it.

The explicit inducing-function map
$f(g)\mapsto\kappa_w(\det g)f(g)$ identifies the twist of the
normalized induction with the normalized induction of
$(\omega_i\kappa_w)|\,\cdot\,|^{(n-1)/2}$.
This proves the actual intertwiner in (HTS.6).  The weight-zero
Artin identity says $\mu(s)=\iota^{-1}\kappa_w(a(s))$.
Multiply it by (HPT.3) to get the last assertion of (HTS.6),
including smoothness and continuity.  Equivalently the finite
part becomes $\phi_i\mu|_{I_w}^{-1}$: its kernel contains the
intersection of the two open kernels, and the product of this
finite part with $\psi_i\mu$ has the same algebraic expression.
This directly checks every field of local algebraicity and
does not require a new choice of the algebraic weight.
\end{proof}
% END HTS5

% BEGIN HTS6
\paragraph{HTS6. Extracting an actual automorphic part with all ordinary subquotients.}
\label{claim:HTS6}
For $J\subseteq\ker x$ and $x.\mathrm{Good}(h)$, there are an
admitted $\lambda$ with $h=\mathrm{htW}(\lambda)$ and one actual
part $Y$ arising from a nonzero
$Y\mathrm{Of}(Q,\iota,\tau_0,e,\xi,\sigma)$, such that
$x.D=\det Y.r$ and every selected local component of $\sigma$
is a subquotient of $\mathrm{PS}(\omega_v)$ with
$\mathrm{OrdChar}(\omega_{v,i},h_i,x.\psi_{v,i})$.
The same $\sigma$ has the full split $\mathrm{LocGlobT}$
comparison with $Y.r\circ p$, its original infinity weights,
unramifiedness outside $V.S'$, and its original $K_c(1)$ vector.
\begin{proof}
Use the part-dependent predicate which says that $Y$ is obtained
from such a $\sigma,e,\xi$ and has its actual local components
and full split compatibility.  This predicate satisfies the
literal $\mathrm{AutYAlt}$ spanning condition: ADP12 decomposes
each original coefficient function into its $Y\mathrm{Of}$
pieces; discard the zero pieces and apply UPA16 to each nonzero
piece, retaining exactly its same $\sigma,e,\xi$ and local
components.  The integral type has open kernel by ITA12 and
the specified injective coefficient map, as required by ADP12.

The top-ideal branch has no such point.  Otherwise HED1 factors
$x$ through the original Hecke quotient.  AFS8 translates its
weight, and HAW10, with the just proved part-dependent predicate,
gives this actual $Y$, its full determinant equality and all
ordinary subquotients.  Naturality HED2 pulls the equality back
to the identical $x.D$, and the exact IsJ character maps make
the characters precisely $x.\psi$, in their given order.
Finally LTP12, LTP14 and LTP15 apply to this same nonzero
$Y\mathrm{Of}$ to give unramifiedness, infinity weights and
the depth-one fixed vector.  LTP17--LTP18 also give its full
type statement for every Teichm\"uller family.  This explains
why the arbitrary-part version of HAW10 is used: a representation
with a cofinite attachment alone would not supply these data.
\end{proof}
% END HTS6

% BEGIN HTS7
\paragraph{HTS7. The twisted point and every input of the ordinary return.}
\label{claim:HTS7}
For the data in (HTS.1), put $R=Y.r\circ p$ using HTS6 and
$R'=R\mu$ on the same vector space.  Then $x'$ is arithmetic
of weight $h$, $R'$ is continuous semisimple,
$x'.\mathrm{IsRep}(R')$, and the actual twisted automorphic
representation $\sigma'$ of HTS2 satisfies every input of HPT17
for $x',R',\lambda$.
\begin{proof}
HED5 gives continuity and semisimplicity of $R$.  Scalar twisting
preserves invariant subspaces and invariant complements, so
$R'$ is semisimple, and its continuity follows from that of
$R,\mu$.  The last line of (HTS.1) and the finite-part formula
in HTS5 show local algebraicity for $x'.\psi$ with the same
strictly increasing weights.  Thus $x'.\mathrm{Arith}(h)$.

The determinant twist relation is a full polynomial-law identity.
Evaluate it on an arbitrary scalar algebra and a finite sum
$\sum a_g g$.  It replaces each term by
$a_g\beta(g)g$; the matrix under $Y.r$ is therefore
$\sum a_g\beta(g)Y.r(g)$.  Pulling along $p$ gives exactly
the matrix of $R'$ on that sum.  Since $x.D=\det Y.r$,
this proves the full $x'.\mathrm{IsRep}(R')$, including scalar
algebras with nilpotents.

HTS6 gives actual components of $\sigma$ with the original
ordinary subquotients.  Goodness in (HTS.1) and HTS5 promote
them to full principal series, then twist them to the full
principal series and ordered characters in (HTS.6).  The
character identity for $x'$ in (HTS.1) identifies those characters
exactly.  HTS2 preserves the infinity weights.  Outside $V.S'$
the quotient $p$ kills inertia; hence so do $\mu$ and $\delta$,
and HTS2 preserves unramifiedness.  At $v_0$ the stipulated
$I_0\subseteq\ker\beta$ gives the same conclusion and preserves
the original nonzero $K_c(1)$-fixed vector, since this subgroup
lies in the same maximal compact.

We check the guarded comparisons rather than asserting that
twisting preserves an unconditional local statement.  At a
split away-$\ell$ place choose any spherical component of
$\sigma'$, with its actual Satake multiset.  Untwist it to the
corresponding component $\tau$ of $\sigma$ by HTS2.
HTS3's second case gives $\mathrm{Good}(\tau)$ and the
$\mathrm{Gd}$ parameter for every correspondence.  These
activate the away branch of the full $\mathrm{LocGlobT}$ from
HTS6.  HTS4 twists that comparison, and GUA14 converts it to
the original spherical Frobenius law.  This proves the entire
$\mathrm{SphEq}$ for $\sigma',R'$.

At a type place choose any actual component $\tau'$ of
$\sigma'$ and its untwist $\tau$.  LTP17--LTP18, applied to
the original nonzero part, give the inertia type of
$\mathrm{rc}(\mathrm{algNorm}\tau)$ for every correspondence
and every Teichm\"uller family.  The type characters are
restrictions of the global characters $t_i\circ p$, so HTS3's
first case gives exactly $\mathrm{Good}(\tau)$ and
$\mathrm{Gd}$ there.  The full original local comparison
from HTS6 and HTS4 therefore yield
$\mathrm{IsWDOf}(\iota,R',\mathrm{rc}(\mathrm{algNorm}\tau'))$
for every such component and correspondence.  This is the
whole $\mathrm{StrongEquivalenceHida}$, with no omitted guard.
All inputs of HPT17 have now been verified on the same objects.
\end{proof}
% END HTS7

% BEGIN HTS8
\paragraph{HTS8. The literal original TwistOk endpoint.}
\label{claim:HTS8}
For the original data in every positive rank, the same ideal
$J$ satisfies $B.\mathrm{TwistOk}(J)$.
\begin{proof}
Take arbitrary $x,x',h,\beta$ satisfying all the original
hypotheses (HTS.1), including the type-inertia and depth-one
conditions.  HTS1--HTS7 construct the actual twisted automorphic
representation and verify each hypothesis of HPT17 for $x'$.
Apply that all-rank return theorem to conclude
$J\subseteq\ker x'$.  The argument uses $\ell$ odd only to
write $\ell^k=2N+1$ and take its unique square-root power;
it includes $k=0$ and $\ell=3$.  The given coefficient DVR,
fixed $\iota$, integral type, level and ideal were never changed.
This proves every quantifier of the original \texttt{TwistOk};
it makes no low-rank restriction and uses no new descent premise.
\end{proof}
% END HTS8
